% k03 King's existence theorem on the cube with scalar sources: the proofs (continues the section of k02_sources.tex)
\subsection{The Cauchy argument with a $c$-number}\label{ssec:k02:CAh}

\begin{theorem}[the Cauchy property with the scheme $c$-number]\label{thm:k02:CAh}
Fix a real source $s=(u,h)$ and let $\sigma$ range over $\{s,0\}$. Suppose there are $k_0$ and non-negative sequences $b_k$, $r_k(\sigma)$, $\theta_k(\sigma)$ such that for all $k\ge k_0$:
\begin{itemize}[leftmargin=2em]
\item[\textup{(U)}] for every $\sigma$ and both members $(k',j')$ of the diagonal pair at $k$,
$e^{-b_k}I^{(j')}_{k'}(\sigma)\le Z_{k'}(\sigma)\le e^{b_k}I^{(j')}_{k'}(\sigma)+r_k(\sigma)Z_{k'}(0)$, with $\sum_kb_k<\infty$, $r_k(0)\le\frac12$, $\sum_kr_k(0)<\infty$ and $\sum_ke^{C_c(h)(k+2)}r_k(s)<\infty$, where $|c_k(h)|\le C_c(h)(1+k)$ (\cref{lem:k02:cnum}\textup{(a)});
\item[\textup{(B)}] on the support of $\chi_{j(k)}$, \eqref{eq:k02:B} holds at $\sigma$ with $\sum_k\theta_k(\sigma)<\infty$.
\end{itemize}
Put $\Theta_k\defeq\theta_k(s)+\theta_k(0)$, $\delta_k\defeq(4b_k+\Theta_k+2r_k(0))\,e^{4b_k+\Theta_k+2r_k(0)}$ and $\epsilon_k\defeq(e^{-c_k(h)}+e^{-c_{k+1}(h)})\,r_k(s)$. Then for all $k\ge k_0$:
\begin{enumerate}[label=\textup{(\roman*)},leftmargin=*]
\item $n_k(s)\le\mathsf M(s)\defeq\exp(\sum_{i\ge k_0}\delta_i)\big(n_{k_0}(s)+\sum_{i\ge k_0}\epsilon_i\big)<\infty$;
\item $|n_{k+1}(s)-n_k(s)|\le\delta_k\mathsf M(s)+\epsilon_k$, so $n_k(s)$ converges to a limit $n(s)\in[0,\mathsf M(s)]$;
\item if \textup{(U)} and \textup{(B)} hold at $-s$ as well, $n_k(s)\ge1/\mathsf M(-s)$ and $n(s)\in[1/\mathsf M(-s),\mathsf M(s)]\subset(0,\infty)$.
\end{enumerate}
If $b_k$, $\Theta_k$ decay geometrically and $r_k(\sigma)\le e^{-ck^2+Ck}$, then $|n_k(s)-n(s)|\le C'5^{-\kappa k}$ for some $\kappa>0$.
\end{theorem}

\begin{proof}
Fix $k\ge k_0$, $j\defeq j(k)$, and let primes denote the $(k+1)$-model after $j+1$ steps.

\emph{Step 1 (one model, \textup{(U)} only).} For one member write $Z(\sigma)$, $I(\sigma)$, $b=b_k$, $r(\sigma)=r_k(\sigma)$. From (U) at $\sigma=0$, $(1-r(0))Z(0)\le e^{b}I(0)$ and $Z(0)\ge e^{-b}I(0)$, so $I(0)>0$. With $\varrho(\sigma)\defeq I(\sigma)/I(0)$ and $N(\sigma)=Z(\sigma)/Z(0)$,
\begin{equation}\label{eq:k02:onemodel}
e^{-2b}(1-r(0))\,\varrho(\sigma)\ \le\ N(\sigma)\ \le\ e^{2b}\varrho(\sigma)+r(\sigma):
\end{equation}
the upper bound from $Z(\sigma)\le e^{b}I(\sigma)+r(\sigma)Z(0)$ and $1/Z(0)\le e^{b}/I(0)$, the lower one from $Z(\sigma)\ge e^{-b}I(\sigma)$ and $1/Z(0)\ge e^{-b}(1-r(0))/I(0)$. No lower bound on $Z(\sigma)$ for $\sigma\ne0$ is used: the remainder enters relative to $Z(0)$ only.

\emph{Step 2 (the vacuum constant).} $\varrho(\sigma)=\hat I(\sigma)/\hat I(0)$, since $I=e^{-v}\hat I$ with one $v$ for all $\sigma$; likewise for the primed member.

\emph{Step 3 (across the pair, \textup{(B)} only).} On the support of $\chi_j$, (B) at $\sigma=s$ reads $S'_N(\cdot;s)=S_N(\cdot;s)+c_k(h)-c_{k+1}(h)+E$ with $|E|\le\theta_k(s)$; the integrals $\hat I(s)$, $\hat I'(s)$ run over the same variables with the same $\chi_j$, so $\hat I'(s)\in e^{c_{k+1}(h)-c_k(h)}[e^{-\theta_k(s)},e^{\theta_k(s)}]\,\hat I(s)$. At $\sigma=0$ ($c_k(0)=0$), $\hat I'(0)\in[e^{-\theta_k(0)},e^{\theta_k(0)}]\,\hat I(0)$. With $\varrho\defeq\varrho(s)$ for the $k$-model and $\varrho'$ for the $(k+1)$-model,
\begin{equation}\label{eq:k02:pair}
e^{-c_{k+1}(h)}\varrho'\ \in\ [e^{-\Theta_k},e^{\Theta_k}]\;e^{-c_k(h)}\varrho .
\end{equation}

\emph{Step 4 (two-sided recursion).} Write $n\defeq n_k(s)$, $n'\defeq n_{k+1}(s)$, $r_0\defeq r_k(0)$, $r_s\defeq r_k(s)$, $\Theta\defeq\Theta_k$. Apply \eqref{eq:k02:onemodel} to the primed member, multiply by $e^{-c_{k+1}(h)}$, use \eqref{eq:k02:pair}, and then \eqref{eq:k02:onemodel} for the $k$-model in the forms $\varrho\le e^{2b}(1-r_0)^{-1}N(s)$ and $\varrho\ge e^{-2b}(N(s)-r_s)$:
\begin{align*}
n'&\le e^{2b}\,e^{-c_{k+1}}\varrho'+e^{-c_{k+1}}r_s\ \le\ e^{2b+\Theta}e^{-c_k}\varrho+e^{-c_{k+1}}r_s\ \le\ e^{4b+\Theta}(1-r_0)^{-1}\,n+e^{-c_{k+1}}r_s,\\
n'&\ge e^{-2b}(1-r_0)\,e^{-c_{k+1}}\varrho'\ \ge\ e^{-2b-\Theta}(1-r_0)\,e^{-c_k}\varrho\ \ge\ e^{-4b-\Theta}(1-r_0)\big(n-e^{-c_k}r_s\big).
\end{align*}
For $t\le\frac12$, $(1-t)^{-1}\le e^{2t}$ and $1-t\ge e^{-2t}$; with $e^{y}-1\le ye^{y}$ and $1-e^{-y}\le y$ this gives $e^{4b+\Theta}(1-r_0)^{-1}-1\le\delta_k$ and $1-e^{-4b-\Theta}(1-r_0)\le\delta_k$, hence
\begin{equation}\label{eq:k02:rec}
(1-\delta_k)\,n_k(s)-\epsilon_k\ \le\ n_{k+1}(s)\ \le\ (1+\delta_k)\,n_k(s)+\epsilon_k .
\end{equation}

\emph{Step 5 (boundedness and convergence).} By induction from the upper half of \eqref{eq:k02:rec}, $n_k\le\prod_{i=k_0}^{k-1}(1+\delta_i)\big(n_{k_0}+\sum_{i<k}\epsilon_i\big)\le\mathsf M(s)$. Here $\sum_i\delta_i<\infty$ because $b_i$, $\Theta_i$, $r_i(0)$ are summable (so bounded, and $\delta_i\le C(4b_i+\Theta_i+2r_i(0))$), and $\sum_i\epsilon_i\le2\sum_ie^{C_c(h)(i+2)}r_i(s)<\infty$ by (U) and $e^{-c_i(h)}\le e^{C_c(h)(1+i)}$. Then \eqref{eq:k02:rec} gives $|n_{k+1}-n_k|\le\delta_kn_k+\epsilon_k\le\delta_k\mathsf M(s)+\epsilon_k$, summable; so $(n_k(s))$ is Cauchy. This is (i), (ii).

\emph{Step 6 (positivity).} $X(-s)=-X(s)$ and $c_k(-h)=-c_k(h)$, so $n_k(s)n_k(-s)=\langle e^{X}\rangle\langle e^{-X}\rangle\ge1$ (\cref{lem:s02:alg}(e)). Steps 1--5 at $-s$ give $n_k(-s)\le\mathsf M(-s)$, hence $n_k(s)\ge1/\mathsf M(-s)$, and the limit inherits both bounds. The rate follows by summing the bound of (ii) from $k$ on.
\end{proof}

\begin{remark}[why the ratio is not taken]\label{rem:k02:additive}
At $h=0$ and $c_k\equiv0$ this is \cref{thm:k01:CA} with \cref{lem:k01:free} replaced by Step~6. The logarithmic form of \cref{thm:k01:CA} needs $r_k(\sigma)Z(0)/Z(\sigma)$ small, that is a lower bound on $Z_k(s)$ uniform enough in $k$; at $h\ne0$ the only candidates are \cref{lem:s02:alg}(e), which needs the upper bound at $-s$ the theorem is to produce, and Jensen's inequality, which needs a renormalised one-point bound that is not available. The additive comparison \eqref{eq:k02:onemodel} avoids both: the remainder is compared with $Z(0)$, and the lower bound on $n_k(s)$ is produced at the end. Nothing in the proof uses the cube or any relation between $j(k+1)$ and $j(k)+1$.
\end{remark}

\subsection{The negative-source bound}\label{ssec:k02:UB}

\begin{theorem}[\UBm: uniform bound at a negative scalar source]\label{thm:k02:UB}
Assume \BALint, \RENorb, \Kh, and \Itwo, \Ithree. Let $0\le q\in C^\infty_c(\R^3)$ be supported at $\ell^\infty$-distance $>1$ from $\R^3\setminus\CJ$. Then $\sup_kn_k(0,-q)<\infty$. Moreover, for every real $2$-form $u$ with the same support condition, $n_k(u,-q)\le n_k(2u,0)^{1/2}n_k(0,-2q)^{1/2}$, so $\sup_kn_k(u,-q)<\infty$.
\end{theorem}

\begin{proof}
By \cref{prop:k02:NEG} choose $M_2>m$ with $s_-(M_2)>\sup q$. \Cref{prop:k01:H0rel,prop:k02:Uplus,prop:k02:Bh} and \cref{lem:k02:cnum} hold at the point $(e,\lambda,M_2,\lmu)$ as well; $\Delta\defeq\sup_k|\Delta_k(M_2)|<\infty$.

\emph{Step 1.} Choose $\bar w\in(0,\frac12]$ with $2\bar w\sup q<s_W$, and $k_1\ge k_0$ with $u_k(M_2)\ge\sup q$ for $k\ge k_1$ (possible since $\sup q<\liminf_ku_k(M_2)$). Then for $k\ge k_1$ and $w\in[0,\bar w]$, both $q\le u_k(M_2)$ and $2wq\le u_k(M_2)$, so \cref{lem:k02:GM} applies at $q$ and at $2wq$.

\emph{Step 2 (an a priori bound at $M_2$).} For $k\ge k_1$, $w\in[0,\bar w]$: by \cref{lem:k02:GM}, $Y^{(M_2)}_k(0,2wq)\le n_k(0,2wq)$; by convexity of $t\mapsto\ln n_k(0,tq)$ with value $0$ at $t=0$, $n_k(0,2wq)\le\max\{1,n_k(0,2\bar wq)\}$; and by \cref{prop:k02:Wsharp}, $n_k(0,2\bar wq)\le e^{\kappa\norm{2\bar wq}_{1,k}}\le e^{\kappa C_0'\wnorm{2\bar wq}}$. Hence, in the conventions of the mass $M_2$,
\[
n^{(M_2)}_k(0,2wq)=e^{\Delta_k(M_2)\sS_k(2wq)}\,Y^{(M_2)}_k(0,2wq)\ \le\ e^{\Delta\,\sS_k(2q)}\max\{1,e^{\kappa C_0'\wnorm{2\bar wq}}\}\ \eqqcolon\ \bar{\mathsf M},
\]
uniformly in $w$ and $k$ ($\sS_k(2q)\le C_0'\wnorm{2q}$).

\emph{Step 3 (the hypotheses of \cref{lem:k02:POS} at $M_2$).} The family $\{wq:\ 0\le w\le\bar w\}$ is compact. At $\sigma=(0,wq)$: (U) holds by \cref{lem:k02:HLF} at mass $M_2$ with $p=2$, from \eqref{eq:k02:SF} (\cref{prop:k02:Uplus} at $M_2$; $wq\ge0$) and \Rzero\ (\cref{prop:k01:H0rel} at $M_2$), with $b_k=C\eta_k^\sigma$ and, in the conventions of $M_2$, $e^{-c_k(wq)}r_k(wq)=\rho_k^{1/2}\,n^{(M_2)}_k(0,2wq)^{1/2}\le\rho_k^{1/2}\bar{\mathsf M}^{1/2}$ by Step~2, a $w$-independent sequence; since $\rho_k\le e^{-ck^2+C}$ and $e^{c_k(wq)}\le e^{C_c(1+k)}$ with $C_c$ uniform on the family, $\sum_ke^{C_c(k+2)}r_k(wq)<\infty$ uniformly in $w$. At $\sigma=0$, (U) is \eqref{eq:k01:H0rel} with $r_k(0)=\rho_k\le\frac12$ for $k\ge k_0$. (B) at $\sigma\in\{wq,0\}$ is \cref{prop:k02:Bh} at $M_2$ on the compact family, with $w$-independent $\theta_k$.

\emph{Step 4 (positivity at $M_2$).} \Cref{lem:k02:POS} gives $c\defeq\inf_{k\ge k_1}n^{(M_2)}_k(0,q)>0$, hence $Y^{(M_2)}_k(0,q)\ge e^{-\Delta\sS_k(q)}c$ for $k\ge k_1$.

\emph{Step 5 (the mass shift).} By \cref{lem:k02:GM} at $q$, $n_k(0,-q)\le1/Y^{(M_2)}_k(0,q)\le e^{\Delta\sS_k(q)}/c\le e^{\Delta C_0'\wnorm q}/c$ for $k\ge k_1$; the finitely many $k<k_1$ are finite numbers. The mixed bound is Cauchy--Schwarz, $\langle e^{F(u)-\Wick{q}}\rangle\le\langle e^{2F(u)}\rangle^{1/2}\langle e^{-2\Wick{q}}\rangle^{1/2}$, with the $c$-numbers distributed ($e^{c_k(q)}=(e^{c_k(2q)})^{1/2}$); the first factor is bounded by \cref{lem:k01:free}.
\end{proof}

\subsection{Existence on the cube with scalar sources}\label{ssec:k02:W1}

\begin{theorem}[King's existence theorem on the cube with scalar sources, reference scheme]\label{thm:k02:W1}
Assume \BALint, \RENorb, \Kh, and \Itwo, \Ithree\ (so that \cref{prop:k02:Wsharp} holds with $s_W=s_1$). Let $f_1,\dots,f_n$ be real $2$-forms and $h_1,\dots,h_p\ge0$ non-negative real functions, all in $C^\infty_c$ and supported at $\ell^\infty$-distance $>1$ from $\R^3\setminus\CJ$, $J\ge J_0$. For $\zeta=(z,w)\in\C^{n+p}$ put $k_1(\zeta)\defeq\sum_j(\Re w_j)_+h_j$, $k_2(\zeta)\defeq\sum_j(\Re w_j)_-h_j$ (both $\ge0$, smooth, with $\sum_j(\Re w_j)h_j=k_1-k_2$), and
\[
D_J(s_W)\ \defeq\ \big\{\zeta\in\C^{n+p}:\ \sup_x2k_1(\zeta)(x)<s_W\big\}.
\]
Then $D_J(s_W)$ is an open convex tube, containing every $\zeta$ with $\Re w_j\le0$ for all $j$ and every $\zeta$ with $w$ purely imaginary, and
\begin{enumerate}[label=\textup{(\roman*)},leftmargin=*]
\item $n_k\to F_J$ locally uniformly on $D_J(s_W)$, and $F_J$ is holomorphic there;
\item $F_J(\zeta)\in(0,\infty)$ for every real $\zeta\in D_J(s_W)$.
\end{enumerate}
\end{theorem}

\begin{proof}
\emph{The domain.} For each $x$, $\Re w\mapsto2k_1(\zeta)(x)=\sum_j2(\Re w_j)_+h_j(x)$ is convex, since $h_j(x)\ge0$ and $t\mapsto t_+$ is convex; so $\Re w\mapsto\sup_x2k_1(\zeta)(x)$ is a finite convex function on $\R^p$, hence continuous, and $D_J(s_W)$ is the preimage under $\zeta\mapsto\Re w$ of an open convex set: open, convex, connected. Write $D\defeq D_J(s_W)$ and $D_\R\defeq D\cap\R^{n+p}$, an open convex subset of $\R^{n+p}$.

\emph{Step 1 (local boundedness).} Let $Q\subset D$ be compact. $\Re Q$ is a compact subset of the open set $D_\R$, so it is covered by finitely many closed coordinate boxes $B\subset D_\R$. For $\zeta\in Q$ with $\Re\zeta\in B$, $|n_k(\zeta)|\le n_k(\Re\zeta)\le\max_{v\in\vertex(B)}n_k(v)$, by \cref{lem:s02:alg}(d) and the convexity of $\ln n_k$ on $B$. At a vertex $v=(x,y)\in D_\R$ put $k_1\defeq k_1(v)$, $k_2\defeq k_2(v)$; then $X(v)=F(u_x)+\Wick{k_1}-\Wick{k_2}$ and the $c$-number is $e^{-c_k(k_1)+c_k(k_2)}$. Hölder's inequality with exponents $(4,2,4)$, with the $c$-number distributed as $e^{-c_k(k_1)}=(e^{-c_k(2k_1)})^{1/2}$, $e^{c_k(k_2)}=(e^{-c_k(-4k_2)})^{1/4}$, gives
\[
n_k(v)\ \le\ a_k(4x)^{1/4}\;n_k(0,2k_1)^{1/2}\;n_k(0,-4k_2)^{1/4}.
\]
The three factors are bounded uniformly in $k$: the first by \cref{lem:k01:free}; the second by \cref{prop:k02:Wsharp}, since $2k_1$ is a non-negative smooth function with $\sup2k_1<s_W$ ($v\in D_\R$) and the required support; the third by \cref{thm:k02:UB} at $q=4k_2$. Hence $\sup_k\sup_Q|n_k|<\infty$.

\emph{Step 2 (convergence on a real open set).} Choose $\delta>0$ with $4\delta\sup_x\sum_jh_j(x)<s_W$ and put $U\defeq\{\zeta=(x,y)\in\R^{n+p}:\ 0<y_j<\delta\ \forall j\}$, a non-empty open subset of $\R^{n+p}$ contained in $D_\R$. At $\zeta\in U$, $k_2(\zeta)=0$ and $s\defeq(u_x,k_1(\zeta))$ has a non-negative scalar part; the sources $\{k_1(\zeta):\zeta\in\bar U\}$ form a compact family. Apply \cref{thm:k02:CAh}(i),(ii) at $s$. Hypothesis (U) at $\sigma=0$ is \eqref{eq:k01:H0rel} on both members, with $r_k(0)=\rho_k$. Hypothesis (U) at $\sigma=s$ is \cref{lem:k02:HLF} with $p=2$, from \eqref{eq:k02:SF} (\cref{prop:k02:Uplus}) and \Rzero, with $r_k(s)=\rho_k^{1/2}\langle e^{X(2s)}\rangle_k^{1/2}$; here $\langle e^{X(2s)}\rangle_k=e^{c_k(2k_1)}n_k(2u_x,2k_1)\le e^{c_k(2k_1)}a_k(4x)^{1/2}n_k(0,4k_1)^{1/2}$ by Cauchy--Schwarz, bounded by $e^{C_c(1+k)}$ times a constant (\cref{lem:k01:free}; \cref{prop:k02:Wsharp}, as $\sup4k_1<s_W$), so $\sum_ke^{C_c(k+2)}r_k(s)\le C\sum_ke^{2C_c(k+2)}\rho_k^{1/2}<\infty$ since $\rho_k\le e^{-ck^2+C}$. Hypothesis (B) at $\sigma\in\{s,0\}$ is \cref{prop:k02:Bh}. Hence $n_k(\zeta)$ converges for every $\zeta\in U$.

\emph{Step 3 (Vitali).} By Step~1 the family $\{n_k\}$ is locally bounded on the connected domain $D$; by Montel's theorem every subsequence has a further subsequence converging locally uniformly to a holomorphic function; any two such limits agree on $U$ (Step~2), hence on $D$, since a holomorphic function vanishing on a non-empty open subset of $\R^{n+p}$ has all its derivatives equal to zero at a point of it and vanishes on the connected $D$. So the whole sequence converges locally uniformly on $D$ to a holomorphic $F_J$: (i).

\emph{Step 4 (positivity at real points).} On $D_\R$ every $n_k$ is positive and $\ln n_k$ is convex, so $F_J\ge0$ there and $g\defeq\ln F_J:D_\R\to[-\infty,\infty)$ satisfies $g(\lambda a+(1-\lambda)b)\le\lambda g(a)+(1-\lambda)g(b)$ (the limit of the inequalities for $\ln n_k$; all values are $<\infty$). For $\zeta\in U$ the point $-\zeta$ lies in $D_\R$ ($k_1(-\zeta)=0$), both $n_k(\zeta)$ and $n_k(-\zeta)$ converge by (i), and $n_k(\zeta)n_k(-\zeta)\ge1$ (\cref{lem:s02:alg}(e)), so $F_J(\zeta)F_J(-\zeta)\ge1$ and $g>-\infty$ on $U$. If $g(t)=-\infty$ for some $t\in D_\R$, pick $t'\in U$ and, $D_\R$ being open, a point $t''\in D_\R$ beyond $t'$ on the ray from $t$ through $t'$, so that $t'=\lambda t+(1-\lambda)t''$ with $\lambda\in(0,1)$; then $g(t')\le\lambda g(t)+(1-\lambda)g(t'')=-\infty$, a contradiction. Hence $F_J>0$ on $D_\R$: (ii).
\end{proof}

\begin{corollary}[\Wone\ in the reference scheme, where it is consumed]\label{cor:k02:W1}
Under the hypotheses of \cref{thm:k02:W1}:
\begin{enumerate}[label=\textup{(\alph*)},leftmargin=*]
\item For every finite family of real $2$-forms $f_i$ and real functions $h_j$ (of either sign) in $C^\infty_c$ with the support condition, and every real coefficient vector $(x,y)$ with $2\norm{(\sum_jy_jh_j)_+}_\infty<s_W$, the limit $\lim_kZ^{z,K}_{k,J}(x,y)$ exists and lies in $(0,\infty)$. In particular \Wone\ for the reference scheme holds at every real point of the tube $\Dtube$ with $\sigma_+=s_1$ (there $2\norm{k_+}_\infty<\sigma_+=s_W$), and at every real coefficient vector whose scalar source $\sum_jy_jh_j$ is non-positive, with no restriction on its amplitude.
\item For $0\le q$ as in \cref{thm:k02:UB}, the comparison-convention functional satisfies $\sup_kN^z_{k,\CJ}(0,-q)\le\sup_kn_k(0,-q)<\infty$.
\end{enumerate}
\end{corollary}
\begin{proof}
(a) Put $\mathsf k\defeq\sum_jy_jh_j$ and let $\mathsf k_+$ be its pointwise positive part. Choose $\chi\in C^\infty_c$ with $0\le\chi\le1$, $\chi=1$ on a neighbourhood of $\supp\mathsf k$ and $\supp\chi$ at distance $>1$ from $\R^3\setminus\CJ$, and $\delta>0$ with $2(\norm{\mathsf k_+}_\infty+\delta)<s_W$; put $h^+\defeq(\norm{\mathsf k_+}_\infty+\delta)\chi$ and $h^-\defeq h^+-\mathsf k$. Both are smooth, non-negative ($h^+\ge\mathsf k_+\ge\mathsf k$ on $\supp\mathsf k$, and $h^+\ge0=\mathsf k$ off it), compactly supported with the support condition, and $\mathsf k=h^+-h^-$. Then $Z^{z,K}_{k,J}(x,y)$ for the given family equals $n_k(\zeta')$ for the family $(f_1,\dots,f_n;h^+,h^-)$ at $\zeta'=(x;1,-1)$, for which $k_1(\zeta')=h^+$ and $\sup2k_1(\zeta')=2(\norm{\mathsf k_+}_\infty+\delta)<s_W$; so $\zeta'\in D_J(s_W)$ and \cref{thm:k02:W1} applies. (b) $e^{-\Wick{q}}$ is a class-$\cD$ observable; killing the matter on $\TJ$ does not decrease its expectation (Step~3 of \cref{lem:s04:G}, the $\cD$ direction), and the scheme multiplier is common.
\end{proof}

\begin{remark}[other schemes]\label{rem:k02:scheme}
For a scheme $a$, $Z^{a,K}_{k,J}(x,y)=e^{-(a_k-z_k)\sum_jy_j\sS_k(h_j)}Z^{z,K}_{k,J}(x,y)$. At a coefficient vector with $\sum_jy_j\int h_j\ne0$ the limit for $a$ therefore exists if and only if $a_k-z_k$ converges: on the domain of \cref{cor:k02:W1}(a), \Wone\ for $a$ is the convergence of the scheme difference, in agreement with \cref{rem:s02:scheme}, and \cref{cor:k02:W1}(b) holds for $a$ whenever $\sup_k|a_k-z_k|<\infty$.
\end{remark}

\begin{remark}[the cell bound]\label{rem:k02:cell}
\Bcell\ (\cref{hyp:s12:Bcell}) is the boundedness of $N^a_{k,C_{L_c}}(0,-2H'\chi_0)$ on one comparison-convention cell. When the cell is a King cube, that is $\ell_c=2MJ'$ for some $J'$ up to a translation by a lattice vector (\cref{lem:s02:pass}), and $\chi_0$ is supported at $\ell^\infty$-distance $>1$ from the cell's complement, \cref{cor:k02:W1}(b) and \cref{rem:k02:scheme} give \Bcell\ for every scheme with $\sup_k|a_k-z_k|<\infty$, and \cref{cor:k02:W1}(a) gives the convergence as well.
\end{remark}

\begin{remark}[what is and is not claimed]\label{rem:k02:scope}
Beyond \BALint, \RENorb\ and \Kh, the only input of this section is the threshold $s_W$ of \cref{prop:k02:Wsharp}, supplied by \Itwo\ and \Ithree. Nothing is claimed at a non-negative scalar source of amplitude $\ge s_W/2$: there the positive-source bound is the content of \Wtwo\ above the threshold (\cref{rem:s03:W2}), which stays open. The negative scalar sources are unrestricted, because the mass shift of \cref{lem:k02:GM} turns a negative source of the mass-$m$ model into a positive source of a heavier model, where the small-field analysis applies with no threshold on the amplitude, and \cref{prop:k02:NEG} supplies a heavy enough mass for every amplitude.
\end{remark}
